\section{Tamarin Model}
\label{sec:model}

\subsection{Cryptographic Primitives and Equational Theory}

Our Tamarin model is built on the following builtins and functions:
\begin{sloppypar}
\begin{itemize}
  \item \textbf{Symmetric encryption/decryption:} We used predefined functions \texttt{senc}/\texttt{sdec} satisfying \texttt{sdec(senc(m,k),k) = m} for symmetric encryption and decryption. The adversary can only decrypt a ciphertext if it knows the corresponding key.
  \item \textbf{Key derivation:} \texttt{kdf(key,label)} is modelled as an uninterpreted two-argument function; no algebraic identities are assumed, which captures the idealised one-way property of a KDF.
  \item \textbf{Message authentication:} Tamarin does not provide a built-in MAC primitive, so we define \texttt{mac(M,K)} and \texttt{macverify(mac(M,K),M,K) = mactrue}; the equation ensures that the adversary can only produce a valid ICV if it knows the ICK.
  \item \textbf{Monotone counters:} We define \texttt{count(x)} to model counter advancement and use it for both MN and KN progression.
  \item \textbf{Message-number generation:} In the model, we proceed with the assumption that the Key Server was pre-selected. Therefore, both the Server and the Key Server must have previously sent MKAPDU messages and a MN. For this reason, we define a private function called \texttt{genMN(\$Party)}.
\end{itemize}
\end{sloppypar}
\subsection{Protocol Participants and Long-Term State}

\paragraph{CAK setup.} The rule \texttt{Setup\_CAK} generates a fresh CAK ${\mathit{cak}}$ and stores it as a persistent shared fact \texttt{!SharedCAK(\textasciitilde cak)}, available to both parties.

\paragraph{Derived keys.} Both participants independently compute
\[
  \mathit{kek} = \mathtt{kdf}(\mathit{cak}, \texttt{'KEK'}),\quad
  \mathit{ick} = \mathtt{kdf}(\mathit{cak}, \texttt{'ICK'})
\]
at the beginning of each protocol step, since both parties share the same CAK.

\subsection{Initial Session Rules}

\paragraph{Initial\_KeyServer\_Send.} The Key Server:
\begin{enumerate}
  \item Retrieves \texttt{!SharedCAK(\textasciitilde cak)}.
  \item Generates a fresh Key Number ${\mathit{kn}}$ and a fresh nonce ${\mathit{KS\_nonce0}}$.
  \item Derives $\mathit{sak1} = \mathtt{kdf}(\mathit{cak}, \langle \$\mathit{Keyserver}, {\mathit{kn}},{\mathit{KS\_nonce0}}\rangle)$.
  \item Computes $\mathit{c1} = \mathtt{senc}(\mathit{sak1}, \mathit{kek})$.
  \item Computes $\mathit{icv1} = \mathtt{mac}(\langle \$\mathit{Keyserver}, \mathit{mn\_ks1},{\mathit{kn}}, \mathit{c1}\rangle, \mathit{ick})$.
  \item Sends $\langle \$\mathit{Keyserver}, \mathit{mn\_ks1}, {\mathit{kn}}, \mathit{c1}, \mathit{icv1}\rangle$ and stores \texttt{KSSTATE0}.
\end{enumerate}
The action fact \texttt{SentInitialSAK} is recorded for use in the authentication lemmas.

\paragraph{Initial\_Server\_Receive.} The Server:
\begin{enumerate}
  \item Receives the MKAPDU from the network.
  \item Retrieves \texttt{!SharedCAK(\textasciitilde cak)} and verifies $\mathit{icv1}$ using a \texttt{\_restrict} action: \texttt{macverify(icv1, \ldots, ick) = mactrue}.
  \item Restricts the received MN to match the expected counter: $\mathit{mn\_ks1} = \mathtt{count}(\mathit{mn\_ks})$.
  \item Decrypts the payload: $\mathit{sak1} = \mathtt{sdec}(\mathit{c1}, \mathit{kek})$.
  \item Computes and sends its acknowledgement ICV2.
  \item Records \texttt{ServerAcceptedSAK} and stores \texttt{SSTATE0}.
\end{enumerate}

\paragraph{Initial\_KeyServer\_Confirm.} The Key Server receives the Server's ACK, verifies ICV2, checks that the received MN matches the expected counter, and transitions to \texttt{KSSTATE1}, recording \texttt{InitialSessionCompleted}.

\subsection{Rekey Rules}

The rekey phase mirrors the initial phase structurally, but it reads the previous state (\texttt{KSSTATE1}/\texttt{SSTATE0}) to enforce monotonicity. Specifically:

\begin{itemize}
  \item $\mathit{mn\_ks2} = \mathtt{count}(\mathit{mn\_ks1})$ and $\mathit{kn1} = \mathtt{count}({\mathit{kn}})$ ensure both counters advance.
  \item A fresh nonce ${\mathit{KS\_nonce1}}$ ensures SAK2 is independent of SAK1 even for identical KN paths.
  
\end{itemize}

\subsection{Adversary Model and Compromise Rules}

We model a standard Dolev--Yao adversary augmented with two compromise rules:
\begin{itemize}
  \item \texttt{Reveal\_SAK}: exposes an initial SAK from \texttt{KSSTATE0}; records \texttt{RevealSAK}.
  \item \texttt{Reveal\_CAK}: exposes the CAK from \texttt{!SharedCAK}; records \texttt{RevealCAK}.
\end{itemize}
These rules allow the formulation of conditional security properties, for example that secrecy holds unless the CAK is revealed, and they capture the forward-secrecy limitations inherent in a pure pre-shared-key protocol.

The rule \texttt{Malicious\_KeyServer\_Send} models a Key Server that uses the CAK to construct a well-formed MKAPDU carrying an arbitrary fresh SAK ($\widetilde{\mathit{fakeSAK}}$) instead of a properly KDF-derived one. This rule is used to probe the protocol's acceptance mechanism.

Quantum adversaries is not thread for this model because all cryptographic primitives in the prothocol is based on symmetric key cryptography. Therefore, the model does not include any quantum-specific adversary rules.

\subsection{Security Properties and Lemmas}
\label{sec:properties}

We specify twelve security properties as Tamarin lemmas. Table~\ref{tab:lemmas} summarises them. The checked entries show that the protocol is executable and that the main agreement, authentication, secrecy, and ordering properties hold under the stated assumptions. The two existence lemmas are especially important: they confirm that the model admits a valid execution trace, while also exposing the structural acceptance weakness in which a malicious Key Server can inject a non-KDF-derived SAK that the Server accepts.

\begin{table}[t]
\centering
\caption{Security lemmas verified against the MKA model.}
\label{tab:lemmas}
\setlength{\tabcolsep}{3pt}
\begin{tabular}{@{}p{5.6cm}p{3.3cm}p{3.3cm}c@{}}
\toprule
\textbf{Lemma} & \textbf{Type} & \textbf{Property} & \textbf{Res.} \\
\midrule
\texttt{executable} & exists-trace & Executability & \cmark \\
\texttt{initial\_agreement} & all-traces & Initial agreement & \cmark \\
\texttt{rekey\_agreement} & all-traces & Rekey agreement & \cmark \\
\texttt{server\_authentication\_initial} & all-traces & Server aliveness & \cmark \\
\texttt{server\_authentication\_rekey} & all-traces & Rekey aliveness & \cmark \\
\texttt{initial\_sak\_secrecy\_no\_reveal} & all-traces & SAK1 secrecy & \cmark \\
\texttt{rekey\_sak\_secrecy\_no\_reveal} & all-traces & SAK2 secrecy & \cmark \\
\texttt{rekey\_cannot\_start\_before\_initial} & all-traces & Ordering & \cmark \\
\texttt{rekey\_after\_initial\_completion} & all-traces & Ordering & \cmark \\
\texttt{key\_update} & all-traces & SAK1 $\neq$ SAK2 & \cmark \\
\texttt{key\_number\_progress} & all-traces & KN monotonicity & \cmark \\
\texttt{fresh\_rekey\_key} & all-traces & Rekey freshness & \cmark \\
\texttt{server\_only\_accepts\_kdf\_keys} & exists-trace & KDF bypass & \cmark \\
\texttt{accepts\_non\_kdf\_sak} & exists-trace & Malicious SAK & \cmark \\
\bottomrule
\end{tabular}
\end{table}


The following paragraphs summarise the principal Tamarin lemmas defined in our model and their intended security meaning. In Tamarin, a lemma is a property asserted over all possible execution traces of the protocol. The notation \texttt{\#i} and \texttt{\#j} denotes time points on a trace: a fact annotated with \texttt{@i} is observed at time \texttt{\#i}. Hence, expressions such as \texttt{j < i} indicate that one event precedes another. The ellipses in the formulas indicate that some fact arguments are omitted for readability; the complete set of lemmas is available in the public repository~\cite{mka-tamarin-repo}.






\subsection{Structural Weakness: KDF-Bypass Attack}

\begin{definition}[Malicious SAK Acceptance]
There exists a trace in which \texttt{MaliciousSAKSent}($sak$) is observed before \texttt{ServerAcceptedSAK}($sak$), and $sak$ was never output by \texttt{KDFGeneratedSAK}.
\end{definition}

This exists-trace lemma captures the scenario in which a Key Server holding the CAK injects a fresh but non-KDF-derived SAK into a well-formed MKAPDU. The Server, having no independent means of verifying whether the SAK was derived through the prescribed KDF path or was arbitrary, accepts it. This is a \emph{structural} property of the protocol rather than a cryptographic break.